\documentclass{article}
\usepackage[letterpaper, margin=1in]{geometry}
\usepackage[utf8]{inputenc}
\usepackage{amsmath}
\usepackage{circuitikz}
\title{Operational Amplifiers\\  \large Introduction to Electric Circuits \\ Supplemental Instruction}
\author{Cooper McAllister}
\date{September 21 2026}
\begin{document}
\maketitle
An Operational Amplifier is commonly drawn with the following symbol: \\ \\
\begin{circuitikz}[american] \draw
(0, 0) node[op amp, noinv input up, anchor=-] (OA) {};
\draw (OA.+) to[short, -o, f<_=$i {=} 0$] ++(-3, 0) to[open, v=$v^+$,] (-3, -2)
(OA.-) to[short, -o, f<=$i {=} 0$] ++(-2, 0) to[open, v=$v^-$] (-2, -2)
(OA.out) to[short, -o] ++(1, 0) to[open, v=$v_o {=} A(v^+ - v^-)$] ++(0, -2.5);
\draw (-3, -2) to[short] (4, -2);
\draw (0, -2) node[ground] {};
\end{circuitikz} \\ \\
It can be replaced with the following circuit: \\ \\
\begin{circuitikz}[american] \draw
(0, 0) to[short] (5, 0) (3, 0) node[ground] {};
\draw (4, 0) to[cV_=$v_o {=} A(v^+ - v^-)$, invert] (4, 3) to[short, -o] (5, 3)
(1.5, 0) to[open, v<=$v^-$] (1.5, 2) to[short, o-] (3, 2)
(0, 0) to[open, v<=$v^+$] (0, 3) to[short, o-] (3, 3);
\end{circuitikz} \\ \\
An important note is that for the first symbol, the supply power connections to the op amp are not shown, so \textbf{KCL cannot be applied} to a network that includes the op amp: \\ \\
\begin{circuitikz}[american] \draw
(0, 0) node[op amp, noinv input up, anchor=-] (OA) {};
\draw (OA.+) to[short, -o] ++(-3, 0)
(OA.-) to[short, -o] ++(-2, 0)
(OA.out) to[short, -o] ++(1, 0);
\draw[red, dashed] (-1, -1) -- (3, -1) -- (3, 2) -- (-1, 2) -- (-1, -1);
\draw[red] (4, -0.5) node[] {No KCL};
\end{circuitikz} \\ \\

\end{document}




